\section{Preliminaries}

Throughout the section, we let $R$ be a ring, and $f,g\in R[x]$ be two non-zero univariate polynomials of degree $p$ and $q$, respectively. Thus, we can write $f=a_p x^p + a_{p-1} x^{p-1} + \cdots + a_0$ and $g=b_q x^q + b_{q-1} x^{q-1} + \cdots + b_0$, with $a_p, b_q \neq 0$.

\subsection{Resultant}

\begin{definition}[Sylvester matrix]
    The Sylvester matrix associated to $f$ and $g$ is the $(p+q)\times(p+q)$ matrix defined as follows:
    \begin{equation*}
        S(f,g)=
        \begin{pmatrix}
            a_p     &        &        &        & b_q     &        &        &        \\
            a_{p-1} & a_p    &        &        & b_{q-1} & b_q    &        &        \\
            \vdots  & \vdots & \ddots &        & \vdots  & \vdots & \ddots &        \\
            a_0     & a_1    & \cdots & a_p    & b_0     & b_1    & \cdots & b_q    \\
                    & a_0    & \ddots & \vdots &         & b_0    & \ddots & \vdots \\
                    &        & \ddots & a_1    &         &        & \ddots & b_1    \\
                    &        &        & a_0    &         &        &        & b_0
        \end{pmatrix},
    \end{equation*}
    where all the non-specified coefficients are zero.
\end{definition}

In Sylvester's original paper \cite{sylvester1853xviii}, no explicit row or column convention was specified for the Sylvester matrix. Consequently, some sources in the literature (e.g., \cite{basu2006algorithms}) present the Sylvester matrix as the transpose of the form given above. In this report, we adopt a column-oriented convention, under which the Sylvester matrix is the matrix of the linear map
\begin{align*}
    \phi : R[x]_{<q} \oplus R[x]_{<p} & \to R[x]_{<p+q} \\
    (u,v)                             & \mapsto uf + vg
\end{align*}
in the canonical monomial basis $(1,x,x^2,\ldots)$. This leads to the definition of the resultant as follows, as well as the following fundamental property of the resultant.

\begin{definition}[Resultant]
    The resultant of $f$ and $g$, denoted by $\Res(f,g)$, is defined as the determinant of the Sylvester matrix $S(f,g)$.
\end{definition}

\begin{proposition}
    Suppose that $R$ is a domain, and let $\mathbb{K}$ be its fraction field. Then, $\Res(f,g)\neq 0$ if and only if $\gcd(f,g)=1$ in $\mathbb{K}[x]$.
\end{proposition}

\begin{proof}
    As $\Res(f,g)$ is the determinant of the linear map $\phi$, we have $\Res(f,g) = 0\iff \ker(\phi)\neq \{0\}$, i.e., there exist non-zero polynomials $u$ and $v$, with $\deg(u)<q$ and $\deg(v)<p$, such that $uf+vg=0$. We now show that this is equivalent to $\gcd(f,g)\neq 1$.

    $(\impliedby)$: Assume $\gcd(f,g)\neq 1$. Let $h=\gcd(f,g)$, and $f_1,g_1\in R[x]$ such that $f=hf_1$ and $g=hg_1$. Then, we can take $u=g_1$ and $v=-f_1$, so that $uf+vg=g_1hf_1-f_1hg_1=0$.

    $(\implies)$: Assume that there exist non-zero polynomials $u$ and $v$, with $\deg(u)<q$ and $\deg(v)<p$, such that $uf+vg=0$. We can write $u=d u_1$ and $g=d g_1$, where $d=\gcd(u,g)$ and $\gcd(u_1,g_1)=1$. Then, we have $u_1 f=-v g_1$, which implies that $g_1\mid u_1 f$. Since $\gcd(u_1,g_1)=1$, we deduce that $g_1\mid f$. Thus, $g_1$ is a common divisor of $f$ and $g$. Moreover, we have $\deg(u)<q=\deg(g)$, which implies that $\deg(d)+\deg(u_1)<\deg(d)+\deg(g_1)$, so that $\deg(u_1)<\deg(g_1)$. Hence, $g_1$ is not constant, and we conclude that $\gcd(f,g)\neq 1$.
\end{proof}
